% !TEX root = ../main.tex
\chapter{The Central Bank and Its Independence}
\label{ch:central-bank}

The earlier chapters treated the money supply as something the monetary authority simply chooses: a single growth factor $z$ in the inflation model of Chapter~\ref{ch:inflation}, a lever pulled at will to collect seigniorage. But who is this authority, and why should we trust it to pull the lever wisely rather than opportunistically? Behind the clean letter $z$ stands an actual institution---in the United States, the Federal Reserve System---with a deliberate internal architecture and a contested degree of insulation from the elected government.

This chapter steps away from formal modelling to describe that institution. We first map the structure of the Federal Reserve System: the twelve regional Reserve Banks, the Board of Governors, and the Federal Open Market Committee through which monetary policy is actually conducted. We then turn to the central political question of monetary economics---\emph{central bank independence}---asking what independence means, why a society might want its central bank insulated from day-to-day politics, and what it gives up in exchange. The European Central Bank serves as a contrasting design at the end. The recurring theme is that the institutions are engineered to balance two goals in tension: technical competence shielded from short-run political pressure, and democratic accountability.

\section{The Federal Reserve System}

The Federal Reserve System is not a single bank but a federated structure with several distinct organs, each with its own membership and powers. The system comprises:

\begin{itemize}
    \item the twelve regional \emph{Federal Reserve Banks};
    \item the \emph{Board of Governors} of the Federal Reserve System, in Washington, D.C.;
    \item the \emph{Federal Open Market Committee} (FOMC), which conducts open market operations;
    \item the \emph{Federal Advisory Council}; and
    \item roughly $2{,}900$ member commercial banks.
\end{itemize}

The design distributes power deliberately. Regional banks anchor the system in local economic conditions and are partly owned by private commercial banks; the Board of Governors is a public body appointed by the federal government; and the FOMC blends the two. We take each organ in turn.

\subsection{Federal Reserve Banks}

The twelve Federal Reserve Banks are \emph{quasi-public} institutions: each is owned by the private commercial banks in its district that are members of the Federal Reserve System, yet each carries out public functions and is overseen by the Board of Governors. This hybrid ownership---private in form, public in mandate---is the first place the system's balancing of interests appears.

Each Reserve Bank has nine directors. Member banks elect six of them, and the Board of Governors appoints the remaining three. The nine seats are organized into three classes, each designed to represent a different constituency:
\begin{itemize}
    \item \textbf{Class A} (three directors, elected by member banks): professional bankers;
    \item \textbf{Class B} (three directors, elected by member banks): prominent leaders from industry, labor, agriculture, or the consumer sector;
    \item \textbf{Class C} (three directors, appointed by the Board of Governors): persons who may \emph{not} be officers, employees, or stockholders of banks, so as to represent the broader public.
\end{itemize}
The nine directors together appoint the president of their Reserve Bank, subject to approval by the Board of Governors. The Class structure is a microcosm of the whole system's intent: bankers are at the table, but they do not sit alone, and a public-interest bloc is built in by construction.

\subsubsection*{Functions of the Reserve Banks}

The Reserve Banks perform the operational work of central banking within their districts. Their functions include:
\begin{itemize}
    \item clearing checks;
    \item issuing new currency and withdrawing damaged currency from circulation;
    \item administering and making discount loans to banks in their districts;
    \item evaluating proposed bank mergers and applications for banks to expand their activities;
    \item acting as liaisons between the business community and the Federal Reserve System;
    \item examining bank holding companies and state-chartered member banks;
    \item collecting data on local business conditions; and
    \item employing staffs of professional economists to research questions bearing on the conduct of monetary policy.
\end{itemize}

\subsubsection*{The Special Role of the Federal Reserve Bank of New York}

Among the twelve, the Federal Reserve Bank of New York occupies a distinctive position, for four reasons.

\begin{enumerate}
    \item Its district contains many of the largest commercial banks in the United States, whose safety and soundness are paramount to the health of the national financial system.
    \item It is actively involved in the bond and foreign exchange markets, where open market operations and intervention are physically carried out.
    \item It is the only Federal Reserve Bank that is a member of the Bank for International Settlements (BIS).
    \item Its president is the only Reserve Bank president who is a \emph{permanent} voting member of the FOMC, serving as the committee's vice-chair.
\end{enumerate}
Together with the chair and vice-chair of the Board of Governors, the president of the New York Fed is therefore one of the three most important officials in the entire system.

\subsubsection*{The Reserve Banks and Monetary Policy}

Although the headline policy decisions are made centrally, the Reserve Banks retain real instruments. Their directors ``establish'' the discount rate (subject to review by the Board of Governors), decide which banks may obtain discount loans, and select one commercial banker from each district to serve on the Federal Advisory Council, which consults with the Board of Governors and supplies information used in conducting policy. Most importantly, five of the twelve Reserve Bank presidents hold a vote on the FOMC at any given time.

\subsubsection*{Member Banks}

Membership in the Federal Reserve System is partly mandatory and partly voluntary. All \emph{national} banks are required to be members; commercial banks chartered by states are not required to join but may choose to. The Depository Institutions Deregulation and Monetary Control Act of 1980 narrowed the practical difference between the two: it subjected \emph{all} depository institutions to the same reserve requirements as member banks and gave all of them access to Federal Reserve facilities.

\subsection{The Board of Governors}

The Board of Governors is the public, central organ of the system, headquartered in Washington, D.C. Its design is engineered to insulate it from short-run political pressure---the structural foundation of the independence debate that occupies the second half of this chapter.

\state{Structure of the Board of Governors}{
\begin{itemize}
    \item \textbf{Seven members}, headquartered in Washington, D.C.
    \item Each is \textbf{appointed by the President} and \textbf{confirmed by the Senate}.
    \item Each serves a single \textbf{14-year, non-renewable term}.
    \item Members are required to come from different Federal Reserve districts.
    \item The \textbf{Chairman} is chosen from among the governors and serves a four-year term.
\end{itemize}
}

The long, non-renewable term is the key feature. Because a governor cannot be reappointed, he or she has no incentive to court the appointing President for a second term; and because the terms are staggered and outlast any single presidency, no one administration can pack the Board quickly. We return to this design in Section~\ref{sec:independence}.

\subsubsection*{Duties of the Board of Governors}

The Board's powers reach across the whole system and the banking sector. It:
\begin{itemize}
    \item votes on the conduct of open market operations (as part of the FOMC);
    \item sets reserve requirements;
    \item controls the discount rate through a ``review and determination'' process;
    \item sets margin requirements (the minimum down payment on securities bought on credit);
    \item sets the salaries of the president and officers of each Reserve Bank and reviews each bank's budget;
    \item approves bank mergers and applications for new activities;
    \item specifies the permissible activities of bank holding companies; and
    \item supervises the activities of foreign banks operating in the United States.
\end{itemize}

\subsubsection*{The Chairman}

The Chairman of the Board of Governors is the public face of the system. The Chairman advises the President on economic policy, testifies before Congress, speaks for the Federal Reserve System to the media, and may represent the United States in negotiations with foreign governments on economic matters.

\subsection{The Research Staff}

The Federal Reserve System is the largest employer of economists not merely in the United States but in the world. The central task of these economists is to follow the incoming economic data from government agencies and private-sector organizations, and to advise policymakers on the likely direction of the economy and the probable impact of monetary policy actions. The seemingly mechanical lever $z$ of Chapter~\ref{ch:inflation} is thus pulled only after a large professional apparatus has weighed in on where it should be set.

\subsection{The Federal Open Market Committee}

Open market operations are the principal tool of U.S. monetary policy, and they are governed by the Federal Open Market Committee.

\state{The FOMC}{
\begin{itemize}
    \item Meets \textbf{eight times a year}.
    \item Consists of the \textbf{seven} members of the Board of Governors, the president of the Federal Reserve Bank of New York, and the presidents of \textbf{four} of the other eleven Reserve Banks (who vote on a rotating basis).
    \item The Chairman of the Board of Governors also chairs the FOMC.
    \item Issues directives to the trading desk at the Federal Reserve Bank of New York, where the operations are executed.
\end{itemize}
}

The arithmetic of the votes is itself a statement of where power lies: the seven governors plus the New York Fed president are permanent voters, so the Washington-appointed bloc commands a structural majority, while the rotating four regional presidents keep the rest of the country represented.

\section{The Independence of Central Banks}
\label{sec:independence}

We now turn to the question that gives this chapter its title. A central bank is \emph{independent} to the extent that it can set monetary policy without being overridden by the elected government. Independence is not all-or-nothing, and a first step is to distinguish two kinds.

\defn{Instrument vs.\ Goal Independence}{
A central bank has \emph{goal independence} if it can set the \emph{objectives} of monetary policy (for example, choosing the target rate of inflation). It has \emph{instrument independence} if, given objectives set elsewhere, it can choose the \emph{tools} used to pursue them (open market operations, the discount rate, reserve requirements) free of political interference. A central bank may have one without the other.
}

\subsection{Gauging the Fed's Independence}

Several structural features bear on how independent the Federal Reserve actually is, and they cut in both directions.

\paragraph{Sources of independence.}
\begin{itemize}
    \item \textbf{Independent revenue.} The Fed funds itself from the interest on its securities holdings and from fees, so it does not depend on annual appropriations from Congress. A body that need not pass the hat is far harder to discipline.
    \item \textbf{Long, staggered, non-renewable terms} on the Board of Governors, as described above, blunt the appointing President's leverage.
\end{itemize}

\paragraph{Sources of dependence.}
\begin{itemize}
    \item \textbf{The Fed's structure is written by Congress and is subject to change at any time.} The Federal Reserve was created by statute and can be reorganized---or its powers curtailed---by statute. Its independence exists at the sufferance of the legislature.
    \item \textbf{Presidential influence.} The President shapes the Fed indirectly through influence on Congress, and directly by appointing the members of the Board of Governors and designating its Chairman. Because the governors' fourteen-year terms are not concurrent with the presidency, this influence accrues only slowly---but over time it is real.
\end{itemize}

\rmkb[Independence is granted, not inherent]{The crucial point is that the Fed's independence is a \emph{policy choice embedded in legislation}, not a constitutional given. Congress could revoke it. The interesting economics, then, is not whether the Fed is independent in some absolute sense, but why a legislature would \emph{choose} to tie its own hands by creating an authority it does not directly control. The next subsection takes up that question.}

\subsection{The Case For and Against Independence}

\subsubsection*{Arguments for independence}

\begin{itemize}
    \item \textbf{Inflationary bias.} Subjecting monetary policy to direct political pressure would impart an inflationary bias: elected officials facing reelection have a standing temptation to stimulate the economy now and let the inflation arrive later.
    \item \textbf{The political business cycle.} Without insulation, policymakers may engineer booms to coincide with elections, generating an artificial cycle in output and inflation timed to the political calendar rather than to economic fundamentals.
    \item \textbf{Deficit accommodation.} A politically controlled central bank could be pressed into financing large budget deficits by creating money to buy government debt---the \emph{accommodation} that, in the inflation model of Chapter~\ref{ch:inflation}, raises the money-growth factor $z$ and with it the inflation tax. Independence is a commitment device against this temptation.
    \item \textbf{The principal--agent problem.} Monetary policy is too important, and too technical, to be run on short electoral horizons. The principal--agent problem---the gap between what is good for the public and what is good for the official---is arguably worse for elected politicians, whose horizon ends at the next election, than for insulated technocrats.
\end{itemize}

\subsubsection*{Arguments against independence}

\begin{itemize}
    \item \textbf{Undemocratic.} A powerful body that is not elected exercises enormous influence over the economy, which sits uneasily with democratic principles.
    \item \textbf{Unaccountable.} If the central bank performs poorly, there is no straightforward electoral mechanism to remove those responsible.
    \item \textbf{Coordination.} Independence makes it harder to coordinate fiscal and monetary policy, which may pull in opposite directions when they ought to act together.
    \item \textbf{Mixed record.} The Fed has not always used its independence successfully---insulation from politics is no guarantee of good policy.
\end{itemize}

\state{The independence trade-off}{
Independence buys a \emph{credible commitment} to low and stable inflation by removing monetary policy from the short electoral horizon---it is, in effect, a device for tying the government's hands against the inflationary bias of Chapter~\ref{ch:inflation}. The price is a loss of direct democratic accountability and of easy fiscal--monetary coordination. The institutional design surveyed in Section~\ref{sec:independence} is the attempt to purchase the first without paying too much of the second.
}

\subsection{The European Central Bank}

The European Central Bank (ECB) offers an instructive contrast in institutional design. Like the Federal Reserve it is a federated system, but the balance between the center and the national components is struck differently.

\paragraph{Distinctive features.}
\begin{itemize}
    \item The \emph{National Central Banks} control their own budgets and, collectively, the budget of the ECB itself---decentralization runs deeper than in the U.S. system.
    \item Monetary operations are not centralized; they are carried out by the national central banks rather than from a single desk.
    \item The ECB does not supervise and regulate financial institutions, a function that in the United States the Fed performs prominently.
\end{itemize}

\paragraph{The Governing Council.}
The ECB's policy is set by its Governing Council, the analogue of the FOMC.
\begin{itemize}
    \item It holds monthly meetings at the ECB in Frankfurt, Germany.
    \item It comprises the heads of the National Central Banks together with the six members of the Executive Board.
    \item It operates by consensus; the ECB announces the target rate and takes questions from the media.
    \item To keep the council at a workable size as new countries join the euro area, voting proceeds on a system of rotation.
\end{itemize}

\paragraph{Independence.}
By statutory design the ECB is among the most independent central banks in the world.
\begin{itemize}
    \item Members of the Executive Board serve long terms.
    \item It determines its own budget.
    \item Its overriding mandate is \emph{price stability}, which makes it \emph{less goal-independent} than the Fed: the goal is fixed for it by treaty, so its independence is largely about instruments rather than objectives.
    \item Most strikingly, its charter cannot be changed by ordinary legislation; it can be amended only by revising the Maastricht Treaty---a far higher bar than the statute that governs the Fed.
\end{itemize}

\rmkb[Two designs, one trade-off]{The Fed and the ECB sit at different points on the goal--instrument spectrum. The ECB has a single, treaty-fixed objective (price stability) but exceptionally strong protection of its instruments, while the Fed enjoys greater goal latitude---its mandate also embraces maximum employment---but rests on a charter that Congress can rewrite. Both are answers to the same underlying problem: how to make a credible commitment to sound money while remaining, in the end, answerable to the public. With the architecture of the monetary authority in place, Chapter~\ref{ch:money-supply} turns to the mechanics by which that authority actually changes the quantity of money.}
